% !TEX program = xelatex
% 2027 新领军学科能力一试 · 数学 2（预测训练卷）
\documentclass[cjk=noto, seal=first]{mojing-exam}

\examtitle{2027 新领军学科能力测试 · 预测训练卷}
\secret{训练用卷 ★ 非官方试题}

\begin{document}

\begin{paper}[time={建议 3 小时（训练设置）}, full=150]{一试 · 数学 2}

\begin{instructions}
  \item 本卷共 6 题，满分 150 分，建议用时 3 小时。官方未公布一试的题量、时长与分值，本卷安排均为训练设置。
  \item 数学 2 对应 2027 年考纲中的概率与统计（含组合）、应用数学（数值分析、数值代数、最优化）。
  \item 解答应写出完整的推理、证明与计算过程；只写结果而无论证的，不能得到相应分数。
\end{instructions}

\begin{gradingnotes}
  \item 本评分标准为训练用建议标准，不是官方评分细则。分值标在关键步骤之后，表示做到该步骤累计可得的分数。
  \item 考生方法与本解答不同时，只要论证严格完整，参照相应步骤的分值给分。
\end{gradingnotes}

\scoretable{1,2,3,4,5,6}

\begin{questions}

% =============================================================================
\question[25]
设 $X_1,X_2,\dots$ 独立同分布，$P(X_i=1)=P(X_i=-1)=\frac12$，记 $S_k=X_1+X_2+\dots+X_k$。对正整数 $n$，令
\[
  V_n=\max\Bigl\{\operatorname{Var}\bigl(a_1S_1+a_2S_2+\dots+a_nS_n\bigr)\ :\ a_1,\dots,a_n\in\R,\ a_1^2+\dots+a_n^2=1\Bigr\}.
\]
\begin{enumerate}
  \item 求 $V_n$；
  \item 求 $\displaystyle\lim_{n\to\infty}\frac{V_n}{n^2}$。
\end{enumerate}
\answerspace{2cm}
\begin{solution}
答案：(1) $V_n=\dfrac{1}{4\sin^2\frac{\pi}{4n+2}}$；(2) $\dfrac4{\pi^2}$。

(1)~令 $b_j=a_j+a_{j+1}+\dots+a_n$，$b_{n+1}=0$。交换求和次序，$\sum_ka_kS_k=\sum_jb_jX_j$；$X_j$ 独立、方差为 $1$，故方差为
$\sum_jb_j^2$。又 $a_j=b_j-b_{j+1}$，$a\mapsto b$ 可逆，所以
\[
  V_n=\max_{b\ne0}\frac{\sum_{j=1}^nb_j^2}{q(b)},\qquad q(b)=\sum_{j=1}^n(b_j-b_{j+1})^2 .
\]
\score{8}
$q$ 正定，矩阵 $Q$ 为三对角阵：对角元为 $1,2,\dots,2$，次对角元为 $-1$。由谱定理，$V_n=1/\lambda_{\min}(Q)$。\score{12}
试 $b_j=\cos\bigl((j-\frac12)\theta\bigr)$：中间各行 $-b_{j-1}+2b_j-b_{j+1}=(2-2\cos\theta)b_j$ 恒成立；第一行等价于 $b_0=b_1$，
由余弦的偶性自动成立；最后一行等价于 $b_{n+1}=0$，即 $\cos\bigl((n+\frac12)\theta\bigr)=0$，
$\theta_k=\frac{(2k-1)\pi}{2n+1}$（$k=1,\dots,n$）。这 $n$ 个值互不相同，给出 $Q$ 的全部特征值 $4\sin^2\frac{\theta_k}{2}$。\score{20}
最小者为 $k=1$，所以 $V_n=\dfrac1{4\sin^2\frac{\pi}{4n+2}}$。（检验：$V_1=1$，$V_2=\frac{3+\sqrt5}2$。）\score{22}

(2)~$\sin x\sim x$，$V_n\sim\dfrac{(4n+2)^2}{4\pi^2}$，极限为 $\dfrac4{\pi^2}$。\score{25}

\emph{另一路线}　$\operatorname{Cov}(S_i,S_j)=\min(i,j)$，$(\min(i,j))=LL\T$（$L$ 为元素全为 $1$ 的下三角阵），其逆为三对角阵。
$4/\pi^2$ 是布朗运动协方差核 $\min(s,t)$ 在 $[0,1]$ 上的最大特征值，本题是它的离散版本。
\end{solution}

% =============================================================================
\question[25]
一个种群按代繁衍：第 $0$ 代有 $1$ 个个体；每个个体独立地产生 $0,1,2$ 个后代的概率分别为 $\frac14,\frac12,\frac14$，
第 $n$ 代全部个体的后代构成第 $n+1$ 代。记 $p_n$ 为第 $n$ 代仍有个体的概率。
\begin{enumerate}
  \item 求 $\displaystyle\lim_{n\to\infty}np_n$；
  \item 求 $\displaystyle\lim_{n\to\infty}\frac{n(np_n-4)}{\ln n}$。
\end{enumerate}
\answerspace{2cm}
\begin{solution}
答案：(1) $4$；(2) $-4$。

(1)~记 $q_n=1-p_n$。对第 $0$ 代个体的后代数取条件：每个孩子的子孙独立地按同一规律繁衍，“到第 $n+1$ 代灭绝”等价于
“每个孩子的子孙在 $n$ 代内灭绝”，所以 $q_{n+1}=\frac14+\frac12q_n+\frac14q_n^2=\frac{(1+q_n)^2}{4}$，即
\[
  p_{n+1}=p_n-\frac{p_n^2}{4},\qquad p_0=1 .
\]
\score{7}
取倒数：$\dfrac1{p_{n+1}}-\dfrac1{p_n}=\dfrac{1}{4-p_n}\in\bigl(\frac14,\frac13\bigr]$。求和得 $1+\frac n4\le\frac1{p_n}\le1+\frac n3$，
于是 $p_n\to0$，$\frac1{4-p_n}\to\frac14$，由 Stolz 定理 $\frac{1/p_n}{n}\to\frac14$，$\lim np_n=4$。\score{12}

(2)~由 $\dfrac1{4-p}=\dfrac14+\dfrac p{16}+\dfrac{p^2}{16(4-p)}$，
\[
  \frac1{p_n}=1+\frac n4+\frac1{16}\sum_{k=0}^{n-1}p_k+R_n,\qquad 0\le R_n\le\frac1{48}\sum_{k\ge0}p_k^2 .
\]
由 $p_k\le\frac4{k+4}$，$R_n=O(1)$，$\sum_{k<n}p_k=O(\ln n)$，代回得 $\frac1{p_k}\le\frac{k+4}4+O(\ln(k+2))$，
从而 $0\le\frac4{k+4}-p_k=O\bigl(\frac{\ln(k+2)}{(k+4)^2}\bigr)$，其和收敛。\score{19}
所以 $\sum_{k<n}p_k=4\ln n+O(1)$，$\dfrac1{p_n}=\dfrac n4+\dfrac14\ln n+O(1)$，
\[
  np_n=\frac{4}{1+\frac{\ln n}{n}+O(\frac1n)}=4-\frac{4\ln n}{n}+O\Bigl(\frac1n\Bigr),\qquad \frac{n(np_n-4)}{\ln n}\to-4 .
\]
\score{25}

\emph{评注}　数值上比值以 $O(1/\ln n)$ 的速度缓慢趋于 $-4$（$n=10^7$ 时约为 $-4.69$），只能靠证明确定。
\end{solution}

% =============================================================================
\question[25]
某批产品按 $1,2,\dots,N$ 连续编号，$N$ 为未知的正整数。从中依次不放回地抽取 $n$ 件（$2\le n\le N$），每次在剩余产品中
等可能地抽取一件，记抽到的最大编号为 $M$、最小编号为 $m$。
\begin{enumerate}
  \item 求 $N$ 的极大似然估计 $\hat N$，并判断它是否为 $N$ 的无偏估计（即是否对一切 $N\ge n$ 都有 $E\hat N=N$）；
  \item 求常数 $a_n,b_n$，使对一切 $N\ge n$，$T_1=a_nM+b_n$ 都是 $N$ 的无偏估计，并求 $\operatorname{Var}(T_1)$；
  \item 若编号改为 $s+1,s+2,\dots,s+N$，其中非负整数 $s$ 也未知。求常数 $c_n,d_n$，使 $T_2=c_n(M-m)+d_n$ 对一切 $s\ge0$、
        $N\ge n$ 都是 $N$ 的无偏估计，并当 $n<N$ 时求 $\operatorname{Var}(T_2)/\operatorname{Var}(T_1)$。
\end{enumerate}
\answerspace{2cm}
\begin{solution}
答案：(1) $\hat N=M$，$n<N$ 时有偏；(2) $a_n=\frac{n+1}n$，$b_n=-1$，$\operatorname{Var}(T_1)=\frac{(N+1)(N-n)}{n(n+2)}$；
(3) $c_n=\frac{n+1}{n-1}$，$d_n=-1$，比值为 $\dfrac{2n}{n-1}$。

\textbf{关键观察：间隔可交换。} 把抽到的编号从小到大排为 $y_1<\dots<y_n$，令 $G_0=y_1-1$，$G_i=y_{i+1}-y_i-1$（$1\le i\le n-1$），
$G_n=N-y_n$。$(G_0,\dots,G_n)$ 与无序样本一一对应，且恰取遍和为 $N-n$ 的全部非负整数组；无序样本等可能，故 $(G_0,\dots,G_n)$
在这些整数组上均匀分布，其分布在下标置换下不变。\score{6}

(1)~有序样本的概率为 $\frac{1}{N(N-1)\cdots(N-n+1)}$（$N\ge M$ 时），否则为 $0$，在 $N\ge M$ 上严格递减，故 $\hat N=M$。
由可交换性 $E G_i=\frac{N-n}{n+1}$，$EM=N-EG_n=\frac{n(N+1)}{n+1}<N$（$n<N$），所以有偏。\score{10}

(2)~$a_n\frac{n(N+1)}{n+1}+b_n=N$ 对一切 $N\ge n$ 成立，得 $a_n=\frac{n+1}n$，$b_n=-1$。记 $v=\operatorname{Var}(G_i)$，
$c=\operatorname{Cov}(G_i,G_j)$（$i\ne j$）。由 $\sum G_i$ 为常数，$(n+1)v+n(n+1)c=0$，$c=-\frac vn$。$\operatorname{Var}(M)=\operatorname{Var}(G_n)=v$。
用 $P(M\le j)=\binom jn/\binom Nn$ 及曲棍球恒等式可算出 $E[M(M+1)]=\frac{n(N+1)(N+2)}{n+2}$，于是
\[
  v=\operatorname{Var}(M)=\frac{n(N+1)(N-n)}{(n+1)^2(n+2)},\qquad \operatorname{Var}(T_1)=\Bigl(\frac{n+1}{n}\Bigr)^2v=\frac{(N+1)(N-n)}{n(n+2)} .
\]
\score{17}

(3)~极差 $M-m=N-1-G_0-G_n$ 与 $s$ 无关，$E(M-m)=N-1-\frac{2(N-n)}{n+1}=\frac{(n-1)(N+1)}{n+1}$，故 $c_n=\frac{n+1}{n-1}$，$d_n=-1$。\score{20}
$\operatorname{Var}(M-m)=\operatorname{Var}(G_0+G_n)=2v+2c=\frac{2(n-1)}{n}v$，所以
\[
  \operatorname{Var}(T_2)=\Bigl(\frac{n+1}{n-1}\Bigr)^2\frac{2(n-1)}{n}v=\frac{2(N+1)(N-n)}{(n-1)(n+2)},\qquad
  \frac{\operatorname{Var}(T_2)}{\operatorname{Var}(T_1)}=\frac{2n}{n-1}.
\]
\score{25}

\emph{评注}　不知道编号起点，代价是方差变成约 $2$ 倍：两端都要估计。$n=2$ 时比值为 $4$。
\end{solution}

% =============================================================================
\question[25]
设有限群 $G$ 作用在集合 $\Omega$ 上（记为 $g\cdot\omega$），$X$ 是取值于 $\Omega$ 的随机变量（若 $\Omega$ 不是有限集，约定下文涉及的映射都可测）。
考虑原假设
\[
  H_0:\ \text{对每个}\ g\in G,\ g\cdot X\ \text{与}\ X\ \text{同分布}.
\]
取定函数 $T\colon\Omega\to\R$，对 $\omega\in\Omega$ 令 $p(\omega)=\dfrac1{\abs G}\#\bigl\{g\in G:\ T(g\cdot\omega)\ge T(\omega)\bigr\}$，
检验规则为：当 $p(X)\le\alpha$ 时拒绝 $H_0$。
\begin{enumerate}
  \item 证明：若 $H_0$ 成立，则对任意 $\alpha\in[0,1]$，有 $P\bigl(p(X)\le\alpha\bigr)\le\alpha$；
  \item 设 $S$ 是 $\Omega$ 上若干个双射组成的有限集合，含恒等映射但不一定构成群。把 $H_0$ 改为“对每个 $s\in S$，$s(X)$ 与 $X$ 同分布”，
        并把 $p$ 的定义中的 $G$ 换成 $S$。(1) 的结论是否对所有这样的 $S$ 仍然成立？证明你的结论；
  \item 某配对试验得到 $6$ 个差值 $d=(3,-1,4,-2,5,6)$。在原假设“各差值相互独立，且每个差值的分布关于 $0$ 对称”下，取
        $G=\{\pm1\}^6$ 按分量变号作用在 $\R^6$ 上，$T(d)=d_1+\dots+d_6$。求 $p(d)$，并判断在显著性水平 $\alpha=0.05$ 下是否拒绝原假设。
\end{enumerate}
\answerspace{2cm}
\begin{solution}
答案：(2) 不一定成立；(3) $p(d)=\frac5{64}\approx0.078$，不拒绝。

(1)~固定 $\omega$，令 $t_h=T(h\cdot\omega)$。由 $h\cdot(g\cdot\omega)=(hg)\cdot\omega$ 及代换 $k=hg$（$h\mapsto hg$ 是 $G$ 上的双射），
\[
  p(g\cdot\omega)=\frac1{\abs G}\#\{k\in G:\ t_k\ge t_g\}.
\]
即 $p(g\cdot\omega)$ 与 $p(\omega)$ 在同一组数 $\{t_k\}$ 中比较，这里用到群的封闭性。\score{8}
\textbf{计数引理}：$J=\{g:\#\{k:t_k\ge t_g\}\le\alpha\abs G\}$ 非空时，取 $g^*\in J$ 使 $t_{g^*}$ 最小，则
$J\subseteq\{k:t_k\ge t_{g^*}\}$，$\abs J\le\alpha\abs G$（含并列值的情形）。故 $N_\alpha(\omega):=\#\{g:p(g\cdot\omega)\le\alpha\}\le\alpha\abs G$。\score{13}
在 $H_0$ 下 $P(p(g\cdot X)\le\alpha)=P(p(X)\le\alpha)$，对 $g$ 平均：
$P(p(X)\le\alpha)=\frac1{\abs G}\sum_gP(p(g\cdot X)\le\alpha)=\frac{E N_\alpha(X)}{\abs G}\le\alpha$。\score{16}

(2)~不一定。取 $\Omega=\{1,2,3\}$，$X$ 均匀分布，$\sigma\colon1\mapsto2\mapsto3\mapsto1$，$S=\{\mathrm{id},\sigma\}$（$\sigma^2\notin S$），
$\sigma(X)$ 仍均匀。取 $T(1)=3$，$T(2)=2$，$T(3)=1$，则 $p(1)=p(2)=\frac12$，$p(3)=1$，$P(p(X)\le\frac12)=\frac23>\frac12$。
原因：$S\cdot1=\{1,2\}$ 而 $S\cdot2=\{2,3\}$，参照集合随观测值改变。\score{21}

(3)~原假设下对任意 $\varepsilon\in\{\pm1\}^6$，$(\varepsilon_id_i)$ 与 $(d_i)$ 同分布，满足 (1) 的条件。$T(d)=15$；$\abs d$ 的分量恰为
$1,\dots,6$，变号后 $T=21-2s$，$s$ 为取负号的绝对值之和。$T\ge15\iff s\le3$，对应 $\varnothing,\{1\},\{2\},\{3\},\{1,2\}$
共 $5$ 种，$p(d)=\frac5{64}>0.05$，不拒绝。\score{25}
\end{solution}

% =============================================================================
\question[20]
对以 $2\pi$ 为周期的连续函数 $f$ 和正整数 $N$，记
\[
  I(f)=\int_0^{2\pi}f(x)\,\dd x,\qquad T_N(f)=\frac{2\pi}N\sum_{k=0}^{N-1}f\Bigl(\frac{2k\pi}N\Bigr).
\]
对周期函数，$T_N(f)$ 就是步长为 $\frac{2\pi}N$ 的复合梯形公式；利用 $f(0)=f(2\pi)$，步长为 $\frac\pi N$（共 $2N$ 个小区间）的
复合 Simpson 公式可以写成 $S_N(f)=\frac{4T_{2N}(f)-T_N(f)}3$。令 $f(x)=\dfrac{1}{5-4\cos x}$。
\begin{enumerate}
  \item 求 $I(f)$ 与 $T_N(f)$ 的精确值；
  \item 求所有正整数 $N$，使 $\abs{S_N(f)-I(f)}>\abs{T_{2N}(f)-I(f)}$。
\end{enumerate}
\answerspace{2cm}
\begin{solution}
答案：(1) $I(f)=\frac{2\pi}3$，$T_N(f)=\frac{2\pi}3\cdot\frac{2^N+1}{2^N-1}$；(2) 全体 $N\ge3$。

(1)~令 $z=\ee^{\ii x}$，则 $5-4\cos x=(2-z)(2-\bar z)$。由等比级数，
$f(x)=\sum_{j,l\ge0}\frac{z^j\bar z^l}{2^{j+l+2}}=\sum_{m\in\Z}c_m\ee^{\ii mx}$，$c_m=\frac{2^{-\abs m}}3$，级数绝对一致收敛，
故 $I(f)=2\pi c_0=\frac{2\pi}3$。\score{6}
单位根的正交性：$\frac1N\sum_{k=0}^{N-1}\ee^{2\pi\ii mk/N}$ 当 $N\mid m$ 时为 $1$，否则为 $0$。交换求和（绝对收敛）：
\[
  T_N(f)=2\pi\sum_{j\in\Z}c_{jN}=\frac{2\pi}3\Bigl(1+2\sum_{j\ge1}2^{-jN}\Bigr)=\frac{2\pi}{3}\cdot\frac{2^N+1}{2^N-1}.
\]
\score{12}

(2)~记 $E_N=T_N-I=\frac{4\pi}3\cdot\frac1{2^N-1}$，则 $S_N-I=\frac{4E_{2N}-E_N}3=\frac{4\pi}9\cdot\frac{3-2^N}{4^N-1}$，
$T_{2N}-I=\frac{4\pi}3\cdot\frac1{4^N-1}$。所以
$\abs{S_N-I}>\abs{T_{2N}-I}\iff\abs{2^N-3}>3\iff N\ge3$。\score{20}

\emph{评注}　$\frac{S_N-I}{T_N-I}\to-\frac13$：Simpson 只把误差缩小约 $3$ 倍，而 $T_{2N}$ 的误差量级由 $2^{-N}$ 变为 $4^{-N}$。
Simpson 公式是梯形公式的 Richardson 外推，假定误差按 $h^2,h^4,\dots$ 展开；对光滑周期函数这些项全为零，误差来自混叠项 $c_{\pm N}$。
\end{solution}

% =============================================================================
\question[30]
设 $0<m<M$。对特征值都在 $[m,M]$ 内的 $n$ 阶实对称正定矩阵 $A$，考虑 $f(x)=\frac12x\T Ax$（$x\in\R^n$）的精确线搜索最速下降法：
给定 $x_0$，若 $g_k=Ax_k\ne0$，令 $x_{k+1}=x_k-\alpha_kg_k$，其中 $\alpha_k$ 使 $f(x_k-\alpha g_k)$ 关于 $\alpha\in\R$ 取最小值。
\begin{enumerate}
  \item 求最小的常数 $c=c(m,M)$，使对一切这样的 $n$、$A$、$x_0$ 及每个使 $g_k\ne0$ 的 $k$，都有 $f(x_{k+1})\le c\,f(x_k)$；
  \item 设 $n=2$，$A=\diag(m,M)$。求所有非零初始点 $x_0$，使 $f(x_1)=c\,f(x_0)$；并证明对这样的 $x_0$，每个 $k$ 都有 $f(x_{k+1})=c\,f(x_k)$。
\end{enumerate}
\answerspace{2cm}
\begin{solution}
答案：(1) $c=\Bigl(\dfrac{M-m}{M+m}\Bigr)^2$；(2) $x_0=t\,(M,\pm m)$，$t\ne0$。

(1)~\textbf{上界。} 记 $g=g_k$。$f(x_k-\alpha g)=f(x_k)-\alpha g\T g+\frac{\alpha^2}2g\T Ag$，最小点
$\alpha_k=\frac{g\T g}{g\T Ag}$，$f(x_{k+1})=f(x_k)-\frac{(g\T g)^2}{2g\T Ag}$。又 $f(x_k)=\frac12g\T A^{-1}g>0$，故
\[
  \frac{f(x_{k+1})}{f(x_k)}=1-\frac{(g\T g)^2}{(g\T Ag)(g\T A^{-1}g)} .
\]
\score{8}
由谱分解 $A=\sum\lambda_ie_ie_i\T$，$\lambda_i\in[m,M]$，有 $\lambda_i+\frac{mM}{\lambda_i}\le m+M$；写 $g=\sum c_ie_i$ 得
$g\T Ag+mM\,g\T A^{-1}g\le(m+M)g\T g$。再由均值不等式，$(g\T Ag)(g\T A^{-1}g)\le\frac{(m+M)^2}{4mM}(g\T g)^2$，故
$\frac{f(x_{k+1})}{f(x_k)}\le1-\frac{4mM}{(m+M)^2}=\bigl(\frac{M-m}{M+m}\bigr)^2$。\score{17}
\textbf{最优性}由 (2) 的例子给出（$k=0$ 时取等），故 $c=\bigl(\frac{M-m}{M+m}\bigr)^2$。\score{19}

(2)~$x_0=(\xi,\eta)\ne0$，$g_0=(m\xi,M\eta)$。二维时特征值只有 $m,M$，第一个不等式自动取等；所以取等当且仅当均值不等式取等：
\begin{align*}
  g\T Ag=mM\,g\T A^{-1}g&\iff m^3\xi^2+M^3\eta^2=mM(m\xi^2+M\eta^2)\\
  &\iff(M-m)(M^2\eta^2-m^2\xi^2)=0\iff M\abs\eta=m\abs\xi,
\end{align*}
即 $x_0=t(M,\pm m)$，$t\ne0$。\score{25}
此时 $g_0=t\,mM(1,\pm1)$，$\alpha_0=\frac2{m+M}$，$x_1=t\cdot\frac{M-m}{M+m}\,(M,\mp m)$，仍是同一形式且非零；归纳得每一步都取等号，
迭代点在两条直线之间“之”字形折返。\score{30}

\emph{评注}　固定步长 $\frac2{m+M}$ 的梯度下降每步也至少把 $f$ 缩小到 $c$ 倍：在最坏情形下，精确线搜索并不比最优固定步长快。
\end{solution}

\end{questions}
\end{paper}
\end{document}
